
\textbf{Corollaire 6.2.} \uline{Avec les hypothèses de} 6.1, \uline{soit} $R\subset K^{\bullet}(X)$ \uline{un sous}-$\lambda$-\uline{anneau engendré par les classes d'un nombre fini de faisceaux}. \uline{Alors la filtration de} $\lambda$-\uline{anneau de} $R$ \uline{est discrète}, \uline{et} $R\cap\Fil^1K^{\bullet}(X)$ \uline{est nilpotent}. \uline{De plus}, $R$ \uline{est un module de type fini sur} $R\cap H^0(X,\underline{\mathbb Z})$ (V 3.9.1).

Comme $X$ est quasi-compact, on peut trouver un recouvrement fini de $X$ par des ensembles ouverts disjoints, tel que sur chacun des ouverts de ce recouvrement les faisceaux dont les classes engendrent $R$ soient tous de rang constant. L'argument standard permet alors de se ramener au cas où ces faisceaux sont de rang constant sur $X$; $R$ est alors augmenté vers $\mathbb Z$.

Soient $\underline{E}_i$, $i=1,\ldots,k$, les faisceaux dont les classes engendrent $R$, $E_i$ leurs classes, $x_i=E_i-\varepsilon(E_i)$. On peut écrire $R$ comme quotient du $\lambda$-anneau engendré par les $x_i$ soumis aux relations $\gamma^j(x_i)=0$ pour $j>r_i=\varepsilon(E_i)$ c'est-à-dire $\mathbb Z[x_i,\gamma^j(x_i)]_{\substack{1\le i\le k\\2\le j\le r_i}}$ ; par suite, $R$ est noethérien. Mais d'après la proposition 6.1, les $x_i$, $\gamma^j(x_i)$ sont nilpotents. Donc l'idéal d'augmentation de $R$ est nilpotent; par ailleurs, il n'est autre que $R\cap\Fil^1K^{\bullet}(X)$.

Tout élément de $\Fil^n R$ s'écrit comme polynome sans composantes de poids $<n$ par rapport aux $x_i$, $\gamma^j(x_i)$. Soit $m$ le plus grand entier tel qu'il existe $i$ tel que $\gamma^{m-1}(x_i)\ne0$; soit $p$ un entier tel que les puissances $p$-ièmes des $x_i$, $\gamma^j(x_i)$ soient toutes nulles. Pour $n>kp(1+\ldots+m)$, on voit tout de suite que $\Fil^nR=0$; la filtration de $R$ est donc discrète.

La filtration de $R$ étant discrète, il n'y a qu'un nombre fini de monômes non nuls par rapport aux $x_i$, $\gamma^j(x_i)$. Ces monômes forment un système de générateurs de $R$ comme $\mathbb Z$-module, ce qui achève la démonstration.

